\documentclass[11pt]{article}
\usepackage[a4paper,margin=27mm]{geometry}
\usepackage{amsmath,amssymb,amsthm,xcolor,booktabs,array}
\newcommand{\PROVED}{\textcolor{green!40!black}{\textbf{PROVED}}}
\newcommand{\OPEN}{\textcolor{red!70!black}{\textbf{OPEN}}}
\newcommand{\AUDIT}{\textcolor{orange!70!black}{\textbf{AUDIT}}}
\newtheorem{theorem}{Theorem}
\newtheorem{proposition}{Proposition}
\newtheorem{lemma}{Lemma}
\title{The Hilbert completion of the odd Poisson quotient: recovery and audit (v97)}
\date{13 August 2026}

\begin{document}
\maketitle

\section{The completion constructed in the earlier run}

The earlier Gamma--Euler run constructed the positive tempered measure
\[
 \nu_{E\Gamma}:=\mu_{E,\mathrm{crit}}*\mu_{\Gamma,1},
 \qquad
 \mu_{E,\mathrm{crit}}=\sum_{n\ge1}\frac1n\delta_{\log n},
\]
and the Hilbert space
\begin{equation}\label{eq:HEG}
 \mathscr H_{E\Gamma}^{\mathrm{crit}}=L^2(\nu_{E\Gamma}).
\end{equation}
It also constructed the gcd GNS completion and proved its finite incidence
comparison with the critical Euler form.  These constructions are \PROVED.

The group $\mathbb R$ acts unitarily on \eqref{eq:HEG} by
\[
 (\mathsf U_uF)(x)=e^{iux}F(x).
\]
Thus \eqref{eq:HEG} is indeed a positive Hilbert completion carrying a
unitary centered scaling group.  The question is whether it is a completion
of the specific quotient
\[
 \mathcal H_-^0=\mathcal H_-/Z\mathcal H_\cap
\]
and whether its Hilbert trace is the odd nuclear character of row (c).

\section{The missing descent map}

Let $\mathcal S$ denote the common test core and let
\[
 \mathcal F_\nu a:=\widehat a\big|_{\operatorname{supp}\nu_{E\Gamma}}
\]
be the sampling/Fourier map whose quadratic form completes to
\eqref{eq:HEG}.  The evident attempted map is
\begin{equation}\label{eq:attemptedmap}
 [h]\longmapsto \mathcal F_\nu(Z^{-1}h).
\end{equation}

\begin{proposition}[The evident map does not descend --- \PROVED]
The map \eqref{eq:attemptedmap} is not well-defined on
$\mathcal H_-/Z\mathcal H_\cap$.
\end{proposition}

\begin{proof}
Replacing $h$ by $h+Zf$ with $f\in\mathcal H_\cap$ replaces
$Z^{-1}h$ by $Z^{-1}h+f$.  Hence well-definedness would require
\[
 \mathcal F_\nu f=0\qquad(f\in\mathcal H_\cap).
\]
But $\mathcal H_\cap$ is defined by only the two Tate conditions
$f(0)=\mathcal Ff(0)=0$ (equivalently, two Mellin moments).  They do not
force the Fourier transform of $f$ to vanish on the infinite support of
$\nu_{E\Gamma}$.  Choose a compactly supported smooth $f$ satisfying the two
linear conditions but with $\widehat f(x_0)\ne0$ at any prescribed
$x_0\ne0$ in that support; such an $f$ exists because three distinct
continuous linear functionals cannot have the third in the span of the first
two.  Then $\mathcal F_\nu f\ne0$.
\end{proof}

Therefore the positive completion constructed previously is a source
boundary/reference space, not yet the odd Poisson quotient.  Identifying the
two without a new map is \AUDIT.

\section{Character mismatch}

\begin{proposition}[Matrix coefficient is not nuclear character --- \PROVED]
The positive-definite Gamma--Euler kernel is a cyclic matrix coefficient of
the unitary representation on \eqref{eq:HEG}.  It is not the operator trace
of that representation.
\end{proposition}

\begin{proof}
In the spectral model, $\mathsf U(f)$ is multiplication by $\widehat f(x)$
on $L^2(\nu_{E\Gamma})$.  Its cyclic matrix coefficient at the constant
vector is
\[
 \langle\mathsf U(f)1,1\rangle=\int\widehat f\,d\nu_{E\Gamma}.
\]
For a nonatomic component of $\nu_{E\Gamma}$, a nonzero multiplication
operator is not trace class.  Hence the last integral is a vector state, not
an operator trace.  Row (c), on the other hand, is the nuclear supertrace of
the quotient representation.  Equality of the two requires a separate
trace-comparison/intertwining theorem.
\end{proof}

This is the precise reason why positivity of the Euler--Gamma GNS space did
not already prove row (d).

\section{The exact obstruction on an off-line residue orbit}

The following finite-dimensional calculation is an audit test for every
proposed Hilbert completion.

\begin{theorem}[Positive Rosati metric excludes a mirror pair --- \PROVED]
Let an off-line functional-equation orbit be represented by two distinct
characters $\rho$ and $\rho'=1-\bar\rho$.  On its residue space the test
algebra acts by
\[
 \pi(f)=\begin{pmatrix}\widehat f(\rho)&0\\0&\widehat f(\rho')\end{pmatrix},
 \qquad
 \pi(f^\vee)=
 \begin{pmatrix}\overline{\widehat f(\rho')}&0\\
                 0&\overline{\widehat f(\rho)}\end{pmatrix}.
\]
There is no positive definite Gram matrix $G$ satisfying
\[
 \pi(f^\vee)=G^{-1}\pi(f)^*G
\]
for every test function $f$.
\end{theorem}

\begin{proof}
Write $G=(g_{ij})$ with $g_{11}>0$.  The adjoint identity is equivalent to
\[
 \begin{pmatrix}\bar a&0\\0&\bar b\end{pmatrix}G
 =G\begin{pmatrix}\bar b&0\\0&\bar a\end{pmatrix},
 \qquad
 a=\widehat f(\rho),\quad b=\widehat f(\rho').
\]
The $(1,1)$ entry gives $(\bar a-\bar b)g_{11}=0$.  Compactly supported test
functions interpolate independently at the two distinct Mellin characters,
so one can choose $a\ne b$.  This forces $g_{11}=0$, contradicting positive
definiteness.
\end{proof}

Thus a faithful positive Hilbert completion with the row-(c) character does
not merely repackage the quotient: its existence itself rules out every
off-line orbit.  This is why preservation of the nuclear character is the
load-bearing clause.

\section{Status of all earlier candidates}

\begin{center}
\begin{tabular}{>{\raggedright\arraybackslash}p{.30\linewidth}
                >{\centering\arraybackslash}p{.13\linewidth}
                >{\centering\arraybackslash}p{.15\linewidth}
                >{\raggedright\arraybackslash}p{.32\linewidth}}
\toprule
Candidate & Positive Hilbert & Unitary scale & Same odd nuclear character \\
\midrule
Tate-renormalized Witt space & yes & rational dilation & no global Poisson
descent constructed \\
gcd GNS space & yes & $\mathbb Q_+^\times$ & no; different spectral type \\
critical Euler form & yes & reference action & no trace comparison \\
critical Euler--Gamma space \eqref{eq:HEG} & yes & continuous unitary action &
no; the evident quotient map fails \\
residue realization & Krein & exact & yes \\
desired odd completion & \OPEN & \OPEN & \OPEN \\
\bottomrule
\end{tabular}
\end{center}

The residue realization is the only previous construction with the exact
character, and it is Hilbert precisely when the off-line mirror blocks are
absent.  The positive GNS constructions have the opposite virtue: they are
Hilbert from the start, but no proved map identifies their vector state with
the odd nuclear trace.

\section{Exact remaining construction}

The required theorem can now be stated without ambiguity:

\medskip
\noindent\textbf{Faithful Poisson--Witt Hilbert descent (\OPEN).}
Construct a map
\[
 \mathfrak Q:\mathcal H_-^0\longrightarrow\mathscr H_{\rm PW}
\]
with dense image such that
\begin{enumerate}
\item $\mathfrak Q\pi_0(u)=\mathsf U_u\mathfrak Q$ for the centered scaling;
\item $\mathfrak Q\pi_-(f^\vee)=\mathfrak Q\pi_-(f)^*$ on the nuclear core;
\item $\operatorname{Tr}_{\mathscr H_{\rm PW}}\pi_-(h)$ equals the odd
      nuclear trace of row (c), not merely a cyclic matrix coefficient.
\end{enumerate}
Then the even Tate trace vanishes on $\mathcal T^0$ and
\[
 B_{\rm nuc}(f,f)
 =-\|\pi_-(f)^*\|_{\rm HS}^2\le0.
\]

The Witt--Tate adjoint required in clause (2) is already \PROVED in v96.
Clauses (1) and (3), jointly, are not proved by any completion recovered from
v89--v94.

\end{document}
